% !TeX program = xelatex
% Hao Pan - Chinese resume
% Overleaf-ready standalone version

\documentclass[a4paper,10pt]{article}

\usepackage[a4paper,left=13.5mm,right=13.5mm,top=9.5mm,bottom=9.5mm]{geometry}
\usepackage{fontspec}
\usepackage{xeCJK}
\usepackage{xcolor}
\usepackage{enumitem}
\usepackage{titlesec}
\usepackage[hidelinks]{hyperref}
\usepackage{fancyhdr}
\usepackage{fontawesome}

\definecolor{accent}{HTML}{2F6661}
\definecolor{project}{HTML}{355B58}
\definecolor{ink}{HTML}{20262E}
\definecolor{muted}{HTML}{5F6975}
\definecolor{rule}{HTML}{CBD5D2}

\setmainfont[
  BoldFont=texgyreheros-bold.otf,
  ItalicFont=texgyreheros-italic.otf,
  BoldItalicFont=texgyreheros-bolditalic.otf
]{texgyreheros-regular.otf}

% Overleaf 优先使用 Noto；其他环境按顺序回退。
\IfFontExistsTF{Noto Sans CJK SC}{
  \setCJKmainfont[BoldFont={Noto Sans CJK SC Bold}]{Noto Sans CJK SC}
  \setCJKsansfont[BoldFont={Noto Sans CJK SC Bold}]{Noto Sans CJK SC}
  \setCJKmonofont{Noto Sans Mono CJK SC}
}{
  \IfFontExistsTF{Hiragino Sans GB}{
    \setCJKmainfont[BoldFont={Hiragino Sans GB W6}]{Hiragino Sans GB W3}
    \setCJKsansfont[BoldFont={Hiragino Sans GB W6}]{Hiragino Sans GB W3}
    \setCJKmonofont{Hiragino Sans GB W3}
  }{
    \setCJKmainfont[BoldFont=FandolHei-Bold.otf]{FandolHei-Regular.otf}
    \setCJKsansfont[BoldFont=FandolHei-Bold.otf]{FandolHei-Regular.otf}
    \setCJKmonofont{FandolHei-Regular.otf}
  }
}

\XeTeXlinebreaklocale "zh"
\XeTeXlinebreakskip=0pt plus 1pt
\XeTeXgenerateactualtext=1
\hypersetup{
  pdftitle={潘昊 - 数据智能体 / 大语言模型应用工程师简历},
  pdfauthor={潘昊},
  pdfsubject={数据智能体 / 大语言模型应用工程师},
  pdfkeywords={Data Agent, LLM Application, Text-to-SQL, Agent Evaluation, Guardrails}
}

\pagestyle{fancy}
\fancyhf{}
\renewcommand{\headrulewidth}{0pt}
\renewcommand{\footrulewidth}{0pt}
\setlength{\parindent}{0pt}
\setlength{\parskip}{0pt}
\setlength{\emergencystretch}{1.2em}
\raggedbottom
\urlstyle{same}
\AtBeginDocument{\fontsize{9.6}{12.8}\selectfont}

\titleformat{\section}
  {\fontsize{12}{14.2}\selectfont\bfseries\color{accent}}
  {}{0pt}{}
  [\vspace{1pt}{\color{rule}\titlerule}]
\titlespacing*{\section}{0pt}{6pt}{3.5pt}

\setlist[itemize]{
  leftmargin=1.15em,
  label={\color{accent}\textbullet},
  itemsep=1.8pt,
  topsep=1.8pt,
  parsep=0pt,
  partopsep=0pt
}

\newcommand{\resumeEntry}[4]{%
  {\fontsize{10.3}{12.1}\selectfont\bfseries #1}%
  \hfill{\fontsize{9.1}{11.1}\selectfont\color{muted}#2}\par
  {\fontsize{9.2}{11.3}\selectfont\color{muted}#3\hfill#4}\par\vspace{-0.5pt}
}

\newcommand{\resumeProject}[4]{%
  {\fontsize{10.3}{12.1}\selectfont\bfseries #1}%
  \hfill{\fontsize{9.1}{11.1}\selectfont\color{muted}#2}\par
  {\fontsize{9.2}{11.3}\selectfont\color{muted}#3\hfill#4}\par\vspace{-0.5pt}
}

\newcommand{\skillline}[2]{%
  \textbf{#1}\hspace{0.45em}#2\par
}

\newcommand{\roleProject}[2]{%
  \par\addvspace{3.8pt}%
  \noindent\hspace*{0.8em}%
  {\fontsize{10}{12}\selectfont\bfseries\color{project}#1}%
  \if\relax\detokenize{#2}\relax\else
    \hspace{0.55em}{\fontsize{9}{10.4}\selectfont\color{muted}（#2）}%
  \fi
  \par\vspace{0.7pt}
}

\newenvironment{roleItems}
  {\begin{itemize}[leftmargin=1.75em,labelsep=0.48em,itemsep=2pt,topsep=1.3pt,parsep=0pt,partopsep=0pt]}
  {\end{itemize}}

\newcommand{\contactsep}{\enspace{\color{rule}\textbar}\enspace}

\begin{document}
\color{ink}

\begin{center}
  {\fontsize{24}{26}\selectfont\bfseries\color{accent} 潘昊}\par
  \vspace{1pt}
  {\fontsize{10.8}{12.5}\selectfont\bfseries 数据智能体 / 大语言模型应用工程师}\par
  \vspace{2pt}
  \small
  \href{tel:+8618226918182}{\faPhone\ +86 18226918182}
  \contactsep
  \href{mailto:hao036@e.ntu.edu.sg}{\faEnvelope\ \underline{hao036@e.ntu.edu.sg}}
  \contactsep
  \href{https://www.linkedin.com/in/hao-pan-59193937a/}{\faLinkedin\ \underline{LinkedIn}}
  \contactsep
  \href{https://github.com/HaoPan036}{\faGithub\ \underline{GitHub}}
  \contactsep
  \href{https://haopan036.github.io/}{\faGlobe\ \underline{个人网站}}
\end{center}
\vspace{-4pt}

\section{个人简介}
我更关心人工智能功能做出来后能不能真正被业务使用，结果能不能复核。能够先与产品经理和业务方明确问题、指标口径和系统边界，再把重复分析做成输入、步骤和检查规则都清楚的智能体技能；也能把分散的历史材料整理成可追溯、由人工选择证据的案例知识库，并通过失败模式检查和评估验证可靠性。

\section{教育背景}
\resumeEntry
  {南洋理工大学}
  {2025.11 -- 至今（预计 2027.01 毕业）}
  {应用人工智能硕士，计算机与数据科学学院}
  {新加坡}
\resumeEntry
  {宿州学院}
  {2019.09 -- 2023.06}
  {数据科学与大数据技术工学学士；国家奖学金（2022--2023，全国前 0.2\%）}
  {中国}

\section{工作经历}
\resumeEntry
  {Shopee Singapore（虾皮）}
  {2026.03 -- 2026.09}
  {人工智能数据分析实习生，买家与营销商业智能团队}
  {新加坡}
\roleProject{数据分析智能体技能开发与可靠性}{}
\begin{roleItems}
  \item 与产品经理和区域业务方确认分析目标、指标口径和交付要求，把 PDP、OPC、Review、Cart、Item 五类重复的 A/B 分析做成实际投入使用的数据分析智能体技能。设计统一的技能编写模板，固定输入参数、SQL、输出格式和返回前检查，单次处理由 \textbf{30--45 分钟缩短到 2--3 分钟}。
  \item 面对 SQL 被改写、查询超时、结果不完整和线上表变化等问题，采用固定版本 SQL、参数校验、按实验组拆分查询和完整性检查。总结 \textbf{15 类失败模式}；在 SQL 版本和参数一致的评估中，\textbf{18 行结果的 19 个字段}均与人工 SQL 一致。
\end{roleItems}
\roleProject{分析案例知识库与人机协同检索智能体}{内部试点}
\begin{roleItems}
  \item 把分散在跟踪表、Google 文档和表格、个人目录中的历史分析材料整理成可追溯的案例包，并统一纳入 GitLab 管理。审查近半年 \textbf{160 条记录}，完成 5 个真实案例的标准化，覆盖 4 份需求、5 份报告、12 个 SQL 和 9 个 CSV；30 个 Markdown 文件全部通过校验（\textbf{0 个错误、0 个警告}）。
  \item 基于 Hangar 交付可运行的内部试点版本：用户输入新需求后，受约束的路由模块在一个已登记的 Diana Topic、历史案例和补充信息之间选择处理路径，再由固定处理模块返回 3--5 个相似案例和原始证据；人工选定 0--3 个案例后，系统才会继续比较并生成固定格式的需求草稿。
  \item 为避免模型引用错误、过期或未确认的材料，将每次检索绑定到一个确定的 Git 提交版本，并分别记录目录、技能和证据文件的 Blob / Hash 来源。案例、Topic、用户身份和选择记录均由后端校验；再用处理模块白名单、状态持久化和幂等控制处理伪造 ID、刷新、重试和后端重启。历史 SQL 只供查看和引用，不自动执行。
\end{roleItems}

\resumeEntry
  {金肯职业技术学院}
  {2023.03 -- 2025.10}
  {讲师兼人工智能技术负责人}
  {中国}
\begin{itemize}
  \item 设计并讲授 Python、应用机器学习和生成式人工智能课程，把抽象概念拆成实验、考核与项目任务，覆盖 \textbf{300+ 名学生}。
  \item 牵头 STEM 机器人项目，指导团队使用 ROS 拆解室内自主飞行与避障问题，获 \textbf{RAICOM 2024 全国总决赛三等奖}。
\end{itemize}

\section{精选项目}
\resumeProject
  {BI Data Agent Sandbox}
  {2026.07}
  {TypeScript、React、Node / Vercel Serverless、AlaSQL、Vitest}
  {\href{https://data-agent-sandbox.vercel.app/}{\underline{在线演示}} $\cdot$ \href{https://github.com/HaoPan036/Data-Agent-Sandbox}{\underline{GitHub}}}
\begin{itemize}
  \item 为解决自由生成 SQL 难复现、结果难审计的问题，搭建一个完全在浏览器运行的 BI 数据智能体沙盒。系统把支持的业务问题转成结构化参数和经过校验的只读 SQL，并基于合成电商数据输出图表、回答和可追查的执行轨迹。
  \item 加入确定性约束、按版本回归评估、失败模式统计和问题样本复核队列。敏感请求会被拦截，边界案例会进入“待人工复核（needs-review）”状态；演示不需要登录、外部服务或模型 API 密钥。
\end{itemize}

\resumeProject
  {ContextCue}
  {2026.06 -- 2026.07}
  {TypeScript、Next.js、React}
  {\href{https://contex-cue.vercel.app/}{\underline{在线演示}} $\cdot$ \href{https://github.com/HaoPan036/ContexCue}{\underline{GitHub}}}
\begin{itemize}
  \item 针对个人智能体跨应用记忆中容易出现的人物串线、召回过期信息和隐私写入问题，搭建上下文层，把模拟语音指令、聊天和人工智能反馈转成结构化记忆，并在回复中保留来源引用。
  \item 用 MemoryGate 控制人物和任务范围、人工确认、TTL 过期，并阻止原始聊天和音频直接写入记忆；设计 \textbf{16 个静态 V0.1 评估场景}，覆盖上下文拼接、人物串线、隐私、过期、证据引用和语气安全。
\end{itemize}

\vspace{-1pt}
\textbf{其他作品：}\href{https://clinical-bazi-web.vercel.app/}{\underline{Clinical Bazi Web}} $|$ 按明确规则排盘，并提供边界清楚的人工智能辅助解读。

\section{技术能力}
\skillline{语言与数据：}{Python、TypeScript、SQL（Presto / Trino、Hive、MySQL）、PySpark、SQLite、Shell}
\skillline{智能体与大语言模型：}{Google ADK、智能体技能 / 操作手册、Text-to-SQL、RAG、大语言模型评估、约束规则、Langfuse}
\skillline{工程工具：}{React、Next.js、Vitest、pytest、Git、Linux、ROS}
\skillline{荣誉：}{华为 ICT 大赛全球总决赛，云赛道全球三等奖（前 5\%，2022.06）}

\end{document}
